\documentclass[unnumsec,webpdf,modern,large,numbered]{oup-authoring-template}
\usepackage{url}
\usepackage{booktabs}
\usepackage{array}
\usepackage{etoolbox}
\setcitestyle{numbers,square,comma}
\def\societylogo{}
\makeatletter
\patchcmd{\ps@opening}{\hfill{\smash{\societylogo{}\hspace*{12pt}}\color{black!20}\rule{45pt}{55pt}}}{}{}{}
\patchcmd{\ps@opening}{\hfill{\smash{\societylogo{}\hspace*{12pt}}\color{black!20}\rule{45pt}{55pt}}}{}{}{}
\makeatother
\journaltitle{NAR Genomics and Bioinformatics}
\pubyear{2026}
\copyrightyear{2026}
\lastpage{10}
\makeatletter
\def\ps@headings{\let\@oddfoot\@empty\let\@evenfoot\@empty\def\@oddhead{\vbox{\hbox to \textwidth{\fontsize{8bp}{10bp}\selectfont\itshape\@journaltitle, \@pubyear\hfill\bfseries\thepage}\vspace{5pt}\rule{\textwidth}{1pt}}}\def\@evenhead{\vbox{\hbox to \textwidth{\fontsize{8bp}{10bp}\selectfont\bfseries\thepage\hfill\itshape\@journaltitle, \@pubyear}\vspace{5pt}\rule{\textwidth}{1pt}}}}
\makeatother
\pagestyle{headings}
\title[Supplementary Data]{Supplementary Data for Layered representation diagnostics for ultra-short metagenomic reads}
\author[1]{Ruixiang Mei}
\author[1]{Jianhua Huang}
\address[1]{School of Data Science, The Chinese University of Hong Kong, Shenzhen, Shenzhen 518172, Guangdong, China}
\abstract{Supplementary figures and source tables supporting the representation-diagnostic analyses, including baseline, estimator, perturbation, parameter, redundancy and runtime audits.}
\keywords{supplementary data, metagenomics, representation diagnostics}
\begin{document}
\renewcommand{\thefigure}{S\arabic{figure}}
\renewcommand{\thetable}{S\arabic{table}}
\maketitle
\section{Supplementary data overview}
This supplementary file collects the visual audits and source-table preview that support the main representation-diagnostic manuscript. Figures S1--S3 evaluate baseline controls, empirical information proxies and sequencing-error-aware perturbations. Figures S4--S7 examine local mutation fraction, P-channel counterfactuals, MSP bin and weight sensitivity, kNN mutual-information robustness and dimension-matched high-k compressed baselines. Figures S8--S10 summarize P/MSP redundancy, runtime, relation audits and the CAMI II marine anonymous-read stability probe. Table S1 gives a compact preview of the CAMI II marine source table; Table S2 reports all prespecified conditional-contribution contrasts from the seven-group K/P/MSP ablation.
\section{Supplementary Figures}
\begin{figure*}[p]
\centering
\includegraphics[width=0.95\textwidth]{figures/supplementary/supp_fig_s1_baseline_audit.pdf}
\caption{Baseline and mixed-metric audit. Panels compare same-dimension PCA/SVD controls, mixed-representation stability and block-weight settings on matched perturbation cells. These controls test whether the main drift pattern can be attributed to feature count, dimensional reduction or one selected K/P/MSP weighting.}
\label{fig:suppfigures1}
{\small\noindent\textbf{Alt text:} Baseline audit plots comparing full-position matrices, reduced controls and block-weight sensitivity.\par}
\end{figure*}

\begin{figure*}[p]
\centering
\includegraphics[width=0.95\textwidth]{figures/supplementary/supp_fig_s2_mi_audit.pdf}
\caption{Estimator-dependent mutual-information (MI) and conditional-MI audit for local change versus matched nuisance perturbation. Values summarize low-dimensional representation-derived distance features; they are descriptive empirical estimates rather than universal information bounds.}
\label{fig:suppfigures2}
{\small\noindent\textbf{Alt text:} Bar plots summarizing empirical mutual-information and conditional-mutual-information audit values across representation blocks.\par}
\end{figure*}

\begin{figure*}[p]
\centering
\includegraphics[width=0.95\textwidth]{figures/supplementary/supp_fig_s3_error_aware_art.pdf}
\caption{Quality-stratified ART perturbation audit. Paired cosine similarity and standardized L2 drift are summarized across low-, middle- and high-quality simulator-derived read strata to test whether the compact stability pattern depends on pooling all ART errors.}
\label{fig:suppfigures3}
{\small\noindent\textbf{Alt text:} Line plots showing ART perturbation behavior across quality levels for paired cosine and L2 drift.\par}
\end{figure*}

\begin{figure*}[p]
\centering
\includegraphics[width=0.95\textwidth]{figures/supplementary/supp_fig_s4_mutation_fraction_sweep.pdf}
\caption{Local mutation-fraction sweep across compact and full-position representations. Curves report the response to increasing structured local-change fractions relative to matched nuisance perturbation; the sweep tests sensitivity across effect sizes rather than a functional-variant classification task.}
\label{fig:suppfigures4}
{\small\noindent\textbf{Alt text:} Line plot showing local mutation-fraction sensitivity across compact representation families.\par}
\end{figure*}

\begin{figure*}[p]
\centering
\includegraphics[width=0.95\textwidth]{figures/supplementary/supp_fig_s5_p_channel_counterfactual_audit.pdf}
\caption{P-channel counterfactual and short-bin reliability audit. Sequence-linked P features are compared with permuted or random added channels, blockwise drift is decomposed, and finite-bin sampling bounds are shown for the shortest read conditions.}
\label{fig:suppfigures5}
{\small\noindent\textbf{Alt text:} Multi-panel audit of P-channel counterfactual behavior, drift decomposition and short-bin sampling bounds.\par}
\end{figure*}

\begin{figure*}[p]
\centering
\includegraphics[width=0.95\textwidth]{figures/supplementary/supp_fig_s6_msp_bin_gamma_sensitivity.pdf}
\caption{MSP binset and weight-sensitivity audit at 69 and 75 bp. Alternative coarse-to-fine binsets and $\gamma\in\{0,0.25,0.5,1,2\}$ are compared using stability and grouped local-change readout, testing whether the default MSP result depends on one fragile bin or weight setting.}
\label{fig:suppfigures6}
{\small\noindent\textbf{Alt text:} Line plots showing MSP binset and gamma-weight sensitivity for 69 and 75 bp reads and delta-readout.\par}
\end{figure*}

\begin{figure*}[p]
\centering
\includegraphics[width=0.95\textwidth]{figures/supplementary/supp_fig_s7_method_hardening_audit.pdf}
\caption{k-nearest-neighbour MI robustness and dimension-matched high-k compressed baseline audit. Estimator settings and subsamples test the stability of the empirical MI increment; hashed and sparse-random-projection canonical $k=15$ vectors test perturbation stability at the 222-feature CK4P-MSP budget.}
\label{fig:suppfigures7}
{\small\noindent\textbf{Alt text:} Multi-panel audit showing kNN mutual-information robustness and dimension-matched high-k compressed baseline comparisons.\par}
\end{figure*}

\begin{figure*}[p]
\centering
\includegraphics[width=0.95\textwidth]{figures/supplementary/supp_fig_s8_redundancy_runtime_audit.pdf}
\caption{P/MSP contribution, redundancy and runtime audit. Blockwise contribution summaries are shown alongside relation measures and feature-extraction time, separating scientific contribution from computational cost. The relation results support overlapping but non-equivalent P and MSP channels.}
\label{fig:suppfigures8}
{\small\noindent\textbf{Alt text:} Multi-panel audit showing P and MSP contributions, redundancy and runtime costs.\par}
\end{figure*}

\begin{figure*}[p]
\centering
\includegraphics[width=0.95\textwidth]{figures/supplementary/supp_fig_s9_p_msp_relation_audit.pdf}
\caption{P/MSP relation audit across the tested read lengths. Canonical-correlation, rank and row-wise alignment summaries quantify the shared biochemical component and the additional multiscale variation in MSP; they are not tests of statistical independence.}
\label{fig:suppfigures9}
{\small\noindent\textbf{Alt text:} Line and heatmap panels showing CCA and correlation relationships between P and MSP feature layers.\par}
\end{figure*}

\begin{figure*}[p]
\centering
\includegraphics[width=0.95\textwidth]{figures/supplementary/supp_fig_s10_cami2_marine_probe.pdf}
\caption{CAMI II marine anonymous-read stability probe. Paired cosine similarity, standardized L2 drift and nearest-clean retrieval are summarized for 4,000 anonymous source reads truncated to 69, 75 and 100 bp under N masking, substitution and combined perturbations. Read-level taxonomic labels were not reconstructed, so this figure supports external stability only.}
\label{fig:suppfigures10}
{\small\noindent\textbf{Alt text:} Line, heatmap and bar panels showing CAMI II marine anonymous-read stability across length and perturbation settings.\par}
\end{figure*}

\section{Supplementary Tables}
\begin{table*}[p]
\centering
\scriptsize
\caption{Preview of the CAMI II marine subset source table. The full CSV is retained in the data package.}
\label{tab:supps10}
\begin{tabular}{@{}p{0.07\textwidth}p{0.12\textwidth}p{0.14\textwidth}p{0.08\textwidth}p{0.13\textwidth}p{0.11\textwidth}p{0.09\textwidth}@{}}
\toprule
\textbf{length} & \textbf{condition} & \textbf{method} & \textbf{features} & \textbf{paired cosine} & \textbf{L2 drift} & \textbf{top-1} \\
\midrule
69 & N\_3pct & CK4 & 136 & 0.9643 & 0.2651 & 1.0000 \\
69 & N\_3pct & CK4+P & 147 & 0.9944 & 0.1059 & 1.0000 \\
69 & N\_3pct & CK4P-MSP & 222 & 0.9948 & 0.1016 & 1.0000 \\
69 & N\_3pct & CK4P-MSP+SD & 297 & 0.9943 & 0.1061 & 1.0000 \\
69 & N\_3pct & CK5 & 512 & 0.9382 & 0.3491 & 1.0000 \\
69 & N\_3pct & CK5+CK4P-MSP & 734 & 0.9665 & 0.2572 & 1.0000 \\
69 & N\_3pct & CSP & 147 & 0.9941 & 0.1080 & 1.0000 \\
75 & N\_3pct & CK4 & 136 & 0.9683 & 0.2497 & 1.0000 \\
75 & N\_3pct & CK4+P & 147 & 0.9950 & 0.0992 & 1.0000 \\
75 & N\_3pct & CK4P-MSP & 222 & 0.9954 & 0.0951 & 1.0000 \\
75 & N\_3pct & CK4P-MSP+SD & 297 & 0.9950 & 0.0994 & 1.0000 \\
75 & N\_3pct & CK5 & 512 & 0.9441 & 0.3322 & 1.0000 \\
75 & N\_3pct & CK5+CK4P-MSP & 734 & 0.9698 & 0.2444 & 1.0000 \\
75 & N\_3pct & CSP & 147 & 0.9948 & 0.1011 & 1.0000 \\
100 & N\_3pct & CK4 & 136 & 0.9687 & 0.2483 & 1.0000 \\
\bottomrule
\end{tabular}

\end{table*}

\begin{table*}[p]
\centering
\scriptsize
\caption{Prespecified conditional-contribution contrasts from the seven-group K/P/MSP ablation. The full representation is CK4P-MSP. Improvement is signed so that positive values favour the full representation; for L2 drift, the sign is reversed because lower values are preferable. Confidence intervals are percentile intervals from 10,000 paired-cell bootstrap resamples. Two-sided Wilcoxon signed-rank tests were adjusted across all 12 rows by the Benjamini--Hochberg procedure. Global cells are length-by-perturbation pairs; local cells are length-by-local-mode grouped cross-validation summaries.}
\label{tab:supps2}
\resizebox{\textwidth}{!}{\begin{tabular}{@{}llllrrl@{}}
\toprule
\textbf{block contrast} & \textbf{evidence layer} & \textbf{metric} & \textbf{$n$} & \textbf{full} & \textbf{comparator} & \textbf{improvement [95\% CI]; BH $q$} \\
\midrule
K $\mid$ P+MSP & global & paired cosine & 24 & 0.989 & 0.999 & $-0.010$ [$-0.013,-0.008$]; $2.4\times10^{-7}$ \\
K $\mid$ P+MSP & global & L2 drift & 24 & 0.129 & 0.027 & $-0.101$ [$-0.118,-0.086$]; $2.4\times10^{-7}$ \\
K $\mid$ P+MSP & global & retrieval & 24 & 1.000 & 0.942 & $0.058$ [$0.036,0.081$]; $7.3\times10^{-4}$ \\
K $\mid$ P+MSP & local & grouped macro-F1 & 12 & 0.979 & 0.970 & $0.008$ [$-0.0002,0.019$]; 0.507 \\
\addlinespace
P $\mid$ K+MSP & global & paired cosine & 24 & 0.989 & 0.984 & $0.005$ [$0.004,0.006$]; $2.4\times10^{-7}$ \\
P $\mid$ K+MSP & global & L2 drift & 24 & 0.129 & 0.156 & $0.027$ [$0.023,0.031$]; $2.4\times10^{-7}$ \\
P $\mid$ K+MSP & global & retrieval & 24 & 1.000 & 1.000 & $0.000$ [$0.000,0.000$]; 1.000 \\
P $\mid$ K+MSP & local & grouped macro-F1 & 12 & 0.979 & 0.978 & $0.001$ [$-0.0002,0.002$]; 0.525 \\
\addlinespace
MSP $\mid$ K+P & global & paired cosine & 24 & 0.989 & 0.984 & $0.005$ [$0.004,0.006$]; $2.4\times10^{-7}$ \\
MSP $\mid$ K+P & global & L2 drift & 24 & 0.129 & 0.156 & $0.027$ [$0.023,0.032$]; $2.4\times10^{-7}$ \\
MSP $\mid$ K+P & global & retrieval & 24 & 1.000 & 1.000 & $0.000$ [$0.000,0.000$]; 1.000 \\
MSP $\mid$ K+P & local & grouped macro-F1 & 12 & 0.979 & 0.904 & $0.074$ [$0.035,0.117$]; $7.3\times10^{-4}$ \\
\bottomrule
\end{tabular}
}
\end{table*}
\clearpage
\end{document}
